\documentclass[11pt,a4paper]{article}
\usepackage[margin=2.5cm]{geometry}
\usepackage{graphicx}
\usepackage{float}
\usepackage{booktabs}
\usepackage{amsmath}
\usepackage{hyperref}

\title{4-bit Array Multiplier: RTL to GDSII\\ \large DVD Assignment}
\author{Nikhil Venkat Atkuru(2024102039) \and Abhinav Venkata Kota(2024102059)}
\date{\today}

\begin{document}
\maketitle
\tableofcontents
\newpage

% ==================================================
\section{Part A - RTL Design and Verification}
% ==================================================

\subsection{Block Diagram}

\begin{figure}[H]
    \centering
    \includegraphics[width=\linewidth,height=0.6\textheight,keepaspectratio]{Block_Diagram_full.png}
    \caption{Block Diagram of the 4-Bit Multiplier using Basic Gates}
    \label{fig:block_diagram}
\end{figure}

The inputs \texttt{a} and \texttt{b} are first stored in two 4-bit registers
(Reg A and Reg B). The 16 AND gates form the partial products, the three rows
of 4 full adders (12 in total) add them row by row, and the 8 product bits are
stored in the output register (Reg P).


\subsection{Simulation Results}

The testbench applies all 256 input combinations
($a = 0 \ldots 15$, $b = 0 \ldots 15$), waits for the 2-cycle latency, and
compares the output $p$ with the expected value $a \times b$. It prints one
line for every mismatch and a single summary line at the end. The below subparts show the respective outputs whenever there is a pass and a fail case.

\subsubsection{PASS Summary Line}

With the correct full adder ($\mathit{sum} = a \oplus b \oplus c_{in}$), all
256 vectors match the expected product and no mismatch lines are printed:

\begin{verbatim}
PASS: 256 vectors checked, 0 errors
\end{verbatim}

Simulation ends at $t = 5140$\,ns. This is 256 vectors at one vector per
10\,ns clock cycle, plus the reset and latency cycles.

\subsubsection{Error Injection Test (FAIL Line)}

To confirm that the testbench actually catches errors, the sum output in adder file was temporarily changed from
$a \oplus b \oplus c_{in}$ to $a \oplus b$ (the carry-in is ignored). The testbench was then re-run, and the change was reverted afterwards.
This is the summary and the comparision lines of the failed result log :

\begin{verbatim}
MISMATCH: a=3 b=3 p=5 expected=9 (t=1060000)
MISMATCH: a=3 b=6 p=10 expected=18 (t=1120000)
MISMATCH: a=3 b=7 p=9 expected=21 (t=1140000)
...
MISMATCH: a=15 b=15 p=181 expected=225 (t=5140000)
FAIL: 256 vectors checked, 95 errors
\end{verbatim}

95 of the 256 vectors fail. The faulty sum is wrong only when a full adder
receives a carry-in of 1. So every product that needs no carries inside the
array still comes out correct: for example, any case with $a$ or $b$ equal to
0, 1, 2, 4 or 8. Every product that does generate a carry is caught. This
shows that the testbench detects the injected fault. The full log is in
\texttt{report/partA/mult\_fail\_results.log}.

\subsection{Hand-Worked Example}

\begin{figure}[H]
    \centering
    \includegraphics[width=0.85\linewidth,height=0.5\textheight,keepaspectratio]{Example.png}
    \caption{Hand Worked Example}
    \label{fig:hand_example}
\end{figure}

Using the last digit of roll number 2024102059, $A = 9 + 5 = 14$
(\texttt{1110}) and $B = 7$ (\texttt{0111}). The four partial products
(\texttt{1110}, \texttt{1110}, \texttt{1110} and \texttt{0000}), each shifted
one place further left, add up to \texttt{01100010} $= 98 = 14 \times 7$.


\subsection{Waveform with 2-Cycle Latency}

\begin{figure}[H]
    \centering
    \includegraphics[width=\linewidth,height=0.4\textheight,keepaspectratio]{waveform_14x7.png}
    \caption{Waveform for the worked out example with proper labels}
    \label{fig:waveform}
\end{figure}

The waveform shows the example $14 \times 7$. Values are in hexadecimal, so
\texttt{E} $= 14$ and \texttt{62} $= 98$. The input \texttt{b} changes to 7
at 4640\,ns while \texttt{a} is 14. Marker A (4645\,ns) is the first rising
edge, where the inputs are stored, and marker B (4655\,ns) is the second
rising edge, where \texttt{p} becomes 98. This confirms the 2-cycle latency.


% ==================================================
\section{Part B - Logic Synthesis with Yosys}
% ==================================================

\subsection{Synthesis Results}

The RTL was synthesized with Yosys using the SkyWater SKY130 standard cell
library \texttt{sky130\_fd\_sc\_hd} (typical corner, 25\,\textdegree C, 1.80\,V).
The design was flattened into a single module, so the 12 full adders were
merged into one block of logic. The results below are taken from
\texttt{synth/mult\_array\_stat.rpt}.

\begin{table}[h]
\centering
\begin{tabular}{|l|l|}
\hline
\textbf{Metric} & \textbf{Value} \\
\hline
Total number of cells & 68 \\
Number of flip-flops & 16 (\texttt{sky130\_fd\_sc\_hd\_\_dfrtp\_1}) \\
Total cell area ($\mu$m$^2$) & 785.754 \\
Flip-flop area ($\mu$m$^2$) & $16 \times 25.024 = 400.384$ \\
Flip-flop area (\% of total) & 50.96\,\% \\
3 most-used cell types & \texttt{dfrtp\_1} (16), \texttt{nand2\_1} (8), \\
 & \texttt{and2\_0} and \texttt{xnor2\_1} (6 each, tied) \\
Netlist simulation result & \texttt{PASS: 256 vectors checked, 0 errors} \\
\hline
\end{tabular}
\caption{Part B synthesis results}
\end{table}

The 16 flip-flops match the design exactly: 4 for input \texttt{a}, 4 for
input \texttt{b} and 8 for the product \texttt{p}. All of them were mapped to
\texttt{dfrtp\_1}, a flip-flop with an active-low reset, which matches the
asynchronous \texttt{rst\_n} in the RTL.

\subsection{Gate-Level Netlist Simulation}

The synthesized netlist \texttt{synth/mult\_array\_netlist.v} was simulated
with the same testbench used in Part A (\texttt{tb/tb\_mult.v}), together with
the SKY130 cell models. All 256 input combinations were checked again, and the
result was:

\begin{center}
\texttt{PASS: 256 vectors checked, 0 errors}
\end{center}

This shows that synthesis did not change the behaviour of the design: the
gate-level netlist gives the same product as the RTL for every input.

\subsection{Report Questions}

\subsubsection{Q1: Logic Cell Count vs.\ 16 AND Gates + 12 Full Adders}

Synthesis produced 68 cells in total. Removing the 16 flip-flops leaves
$68 - 16 = 52$ logic cells.

This is not simply $16 + 12 = 28$, because a full adder is not available as a
single cell in our result. Each full adder needs a sum output and a carry
output, so the tool builds it from several smaller cells. For example,
\texttt{maj3\_1} produces the carry and \texttt{xor3\_1}/\texttt{xnor3\_1} or
\texttt{xnor2\_1} cells produce the sum.

The count is also not simply AND gates plus full adders, because the tool
optimizes the whole design together:
\begin{itemize}
  \item Some adders have a constant 0 input (the first carry-in of each row,
        and the top bit passed down from row 0), so parts of them reduce to a
        simple XOR or AND.
  \item Some AND gates are merged into larger cells such as \texttt{and4\_1},
        \texttt{a21oi\_1} and \texttt{a22oi\_1}, which do the job of several
        gates in one cell.
  \item Logic that is shared between neighbouring adders is built only once.
\end{itemize}
So the 52 logic cells are the tool's smallest-area way of building the same
function, not a one-to-one copy of the 16 AND gates and 12 full adders.

\subsubsection{Q3: Flip-Flop Share of Total Cell Area}

The 16 flip-flops take $400.384\,\mu$m$^2$ out of $785.754\,\mu$m$^2$, which
is 50.96\,\% of the total cell area.

This is large for two reasons:
\begin{itemize}
  \item A flip-flop is much bigger than a simple gate. One \texttt{dfrtp\_1}
        is $25.024\,\mu$m$^2$, while one \texttt{nand2\_1} is about
        $3.75\,\mu$m$^2$, so one flip-flop is roughly the size of 6--7 NAND
        gates.
  \item The logic in a 4-bit multiplier is very small (only 52 cells), while
        the number of flip-flops is fixed by the design at 16. So the
        flip-flops end up taking about half of the area.
\end{itemize}
For a larger multiplier the logic grows faster than the number of
flip-flops, so this percentage would become smaller.


% ==================================================
\section{Part C - Static Timing Analysis with OpenSTA}
% ==================================================

\subsection{Finding the Maximum Frequency}

Static timing analysis checks every register-to-register path against the
clock period without running a simulation. For a path that launches at one
flip-flop and is captured at the next, the data must arrive and settle before
the capturing clock edge. This gives the setup constraint:
\[
  T \;\geq\; t_{clk\to Q} + t_{logic} + t_{setup}
\]
The \emph{slack} of a path is the margin left over:
\[
  \text{slack} = \underbrace{(T - t_{setup})}_{\text{required time}}
               - \underbrace{(t_{clk\to Q} + t_{logic})}_{\text{arrival time}}
\]
Positive slack means the path meets timing; negative slack means it is too
slow. The worst slack (WNS) belongs to the slowest path, the
\emph{critical path}. Because slack falls one-for-one as $T$ shrinks, the
smallest period that still passes is
\[
  T_{\min} = T - \text{WNS}, \qquad F_{\max} = \frac{1000}{T_{\min}}\ \text{MHz}
\]
All runs use the typical corner (TT, 25\,\textdegree C, 1.80\,V), an ideal
clock (no clock tree yet), and an input and output delay of 1.0\,ns.

\subsubsection{Run at $T = 10$\,ns}

At $T = 10$\,ns the worst setup slack is \textbf{+7.027\,ns}, so the design
meets timing with a large margin. The critical path has:

\begin{itemize}
  \item Arrival time: $t_{clk\to Q} + t_{logic} = 0.410 + 2.432 = 2.842$\,ns
  \item Required time: $T - t_{setup} = 10.000 - 0.131 = 9.869$\,ns
  \item Slack: $9.869 - 2.842 = 7.027$\,ns
\end{itemize}

The reset path (asynchronous reset recovery check) and the input/output paths
all have more slack than this, so the register-to-register path through
the array sets the speed of the design. From this run:
\[
  T_{\min} = 10 - 7.027 = 2.973\ \text{ns} \approx 2.97\ \text{ns}
\]

\subsubsection{Run at $T = T_{\min}$}

Re-running at $T = 2.97$\,ns gives a worst slack of \textbf{$-0.003$\,ns}:

\begin{itemize}
  \item Required time: $2.970 - 0.131 = 2.839$\,ns
  \item Arrival time: $2.842$\,ns (unchanged, since the logic delay does not
        depend on the clock period)
  \item Slack: $2.839 - 2.842 = -0.003$\,ns
\end{itemize}

The slack is essentially zero, which confirms the calculation. The small
negative value appears only because the exact $T_{\min}$ of 2.973\,ns was
rounded down to 2.97\,ns. A period of 2.98\,ns would give a slack of
+0.007\,ns. Note that the arrival time is identical in both runs: changing the
clock period only moves the required time, which is why
$T_{\min} = T - \text{WNS}$ works.

\subsection{Timing Results}

\begin{table}[H]
  \centering
  \begin{tabular}{ll}
    \toprule
    \textbf{Metric} & \textbf{Value} \\
    \midrule
    Worst slack at $T = 10$\,ns & 7.027\,ns \\
    $T_{\min}$ & 2.97\,ns \\
    $F_{\max}$ & 336.70\,MHz \\
    Slack at $T = T_{\min}$ & $-0.003$\,ns ($\approx 0$) \\
    Critical path start point & Input register \texttt{b\_r[0]} \\
    Critical path end point & Output register \texttt{p[7]} \\
    Number of cells on the critical path & 10 (2 flip-flops + 8 logic cells) \\
    Latency ($2 \times T_{\min}$) & 5.94\,ns \\
    Throughput & 336.70 million results/s \\
    \bottomrule
  \end{tabular}
  \caption{Part C static timing results (typical corner)}
  \label{tab:partC_results}
\end{table}

The multiplier has two register stages: inputs are captured at the first
clock edge and the product at the second. The latency is therefore
two clock periods, $2 \times 2.97 = 5.94$\,ns. Because a new pair of inputs
can enter on every clock edge, the throughput is one result per
cycle, which equals $F_{\max}$: 336.70 million results per second.

\subsection{Critical Path}

\subsubsection{Start Point and End Point}

\begin{itemize}
  \item \textbf{Start point:} the input register holding \texttt{b[0]}.
  \item \textbf{End point:} the output register for \texttt{p[7]}, the most
        significant product bit, whose input is the final carry out of
        row~3.
\end{itemize}

This matches the structure of the array. Bit \texttt{b[0]} forms partial
product $pp_0$, which enters the array at the very top (row~0 into row~1).
Bit \texttt{p[7]} is the last carry produced at the bottom-left of the array
(row~3). A signal starting at \texttt{b[0]} therefore has to travel down
through every row and then along the carry chain of the last row. This is
the longest route through the multiplier.

\subsubsection{Cells on the Critical Path}

\begin{table}[H]
  \centering
  \begin{tabular}{clccl}
    \toprule
    \textbf{\#} & \textbf{Cell} & \textbf{Delay (ns)} & \textbf{Arrival (ns)} & \textbf{Role} \\
    \midrule
    1  & \texttt{dfrtp\_1}  & 0.410 & 0.410 & Input flip-flop (clock to Q) \\
    2  & \texttt{and4\_1}   & 0.223 & 0.634 & Partial product generation \\
    3  & \texttt{o21ai\_0}  & 0.195 & 0.828 & First adder logic \\
    4  & \texttt{nand3\_1}  & 0.140 & 0.969 & First adder logic \\
    5  & \texttt{maj3\_1}   & 0.380 & 1.349 & Full-adder carry \\
    6  & \texttt{xnor3\_1}  & 0.393 & 1.741 & Full-adder sum \\
    7  & \texttt{maj3\_1}   & 0.369 & 2.110 & Full-adder carry \\
    8  & \texttt{maj3\_1}   & 0.383 & 2.493 & Full-adder carry \\
    9  & \texttt{maj3\_1}   & 0.349 & 2.842 & Full-adder carry (to \texttt{p[7]}) \\
    10 & \texttt{dfrtp\_1}  & setup 0.131 & --- & Output flip-flop \\
    \bottomrule
  \end{tabular}
  \caption{Cells on the critical path, in order}
  \label{tab:critical_path}
\end{table}

The \texttt{maj3} cell is a 3-input majority gate,
$\text{maj}(a,b,c) = ab + bc + ca$, which is exactly the carry-out of a full
adder. The \texttt{xnor3} cell is a 3-input XNOR, which produces the full-adder
sum (the inversion is absorbed by the surrounding logic). The synthesis tool
mapped most full adders onto these two cells. In total, the combinational
logic takes 2.432\,ns, and the flip-flop overhead
($t_{clk\to Q} + t_{setup}$) takes 0.541\,ns, about 18\% of $T_{\min}$.

\subsubsection{Comparison with the $2N$ Full-Adder Estimate}

For an $N$-bit array multiplier, the longest path is roughly $2N$ full-adder
delays: about $N$ adders going down through the rows, and about $N$ more
rippling along the carry chain of the last row. For $N = 4$ this predicts
8 full-adder delays.

The actual path has 5 full-adder cells at the start. The path is shorter than the estimate for the following reasons:
\begin{enumerate}
  \item \textbf{Constant inputs simplify adders.} The first adder in every
        row has a carry-in of 0, and the top input to row~1 is also 0.
        Synthesis reduces these full adders to half adders or plain gates,
        which removes full-adder stages from the path.
  \item \textbf{Logic restructuring.} The synthesis tool does not keep the
        full adders as separate blocks. It merges the AND gates and the first
        adder stage into a few compact gates (\texttt{and4},
        \texttt{o21ai}, \texttt{nand3}).
  \item \textbf{The carry path is fast.} Most of the path follows carry
        outputs (\texttt{maj3}), so a full adder's carry is passed straight
        on without first going through its slower sum output.
\end{enumerate}

The $2N$ figure is therefore a structural upper estimate. It still captures
the key trend: the delay grows linearly with $N$, because the signal must
cross the rows and then ripple along the last row.

\subsection{Critical Path and the $2N$ Estimate}

The critical path starts at the input register for \texttt{b[0]} and ends at
the output register for \texttt{p[7]}. It contains 10 cells: the launching
flip-flop, eight logic cells and the capturing flip-flop. Its
arrival time is 2.842\,ns, and adding the 0.131\,ns setup time gives
$T_{\min} = 2.97$\,ns and $F_{\max} = 336.7$\,MHz. Only 5 of the logic cells are
full-adder cells, compared with the estimate of $2N = 8$ full-adder delays.
The difference comes from synthesis simplifying the adders that have constant
zero inputs, merging the first stage into compact gates, and routing the path
mostly through the fast majority-gate carry outputs. The estimate is an upper
bound that ignores these optimisations, but the path's shape, down through
the rows and then along the final carry chain, matches the theory.


% ==================================================
\section{Part D - Physical Design: RTL to GDSII}
% ==================================================

\subsection{Flow Setup}

The RTL was taken to a GDSII layout with LibreLane, which runs Yosys,
OpenROAD, Magic, KLayout and Netgen through synthesis, floorplan, placement,
clock tree synthesis, routing, extraction and signoff. The design uses a fixed
$150\,\mu$m $\times$ $150\,\mu$m die (\texttt{FP\_SIZING = "absolute"},
\texttt{DIE\_AREA = [0, 0, 150, 150]}).

The starting clock period was $1.2 \times T_{\min}$ from Part C, rounded up
to the next 0.5\,ns: $1.2 \times 2.97 = 3.564 \rightarrow 4.0$\,ns.

\subsection{Timing Closure: Run History}

Unlike Part C, LibreLane signs off timing at nine process, voltage and
temperature corners, including the slow corner (SS, 100\,\textdegree C,
1.60\,V), with extracted wire parasitics, a real clock tree and a
0.25\,ns clock uncertainty. The worst setup slack across all corners
(\texttt{timing\_\_setup\_\_ws} in \texttt{metrics.json}) is always set by the
slow corner. Following the handout, the period was increased by 0.5\,ns after
every run that failed setup.

\begin{table}[H]
  \centering
  \begin{tabular}{cccccc}
    \toprule
    \textbf{Run} & \textbf{Period (ns)} & \textbf{Worst setup slack (ns)} &
    \textbf{Slack at TT (ns)} & \textbf{DRC / LVS} & \textbf{Result} \\
    \midrule
    1 & 4.0 & $-1.254$ & 1.167 & 0 / 0 & Fail (setup) \\
    2 & 4.5 & $-0.754$ & 1.667 & 0 / 0 & Fail (setup) \\
    3 & 5.0 & $-0.254$ & 2.167 & 0 / 0 & Fail (setup) \\
    4 & 5.5 & $-0.009$ & 2.595 & 0 / 0 & Fail (setup) \\
    5 & 6.0 & $-0.001$ & 2.929 & 0 / 0 & Fail (setup) \\
    6 & 6.0\,$^{*}$ & $+0.043$ & 2.966 & 0 / 0 & Pass (not used) \\
    \textbf{7} & \textbf{6.5} & $\mathbf{+0.499}$ & \textbf{3.429} & \textbf{0 / 0} & \textbf{Pass (final)} \\
    \bottomrule
  \end{tabular}
  \caption{LibreLane runs. $^{*}$Run 6 added the optional post-global-routing
  timing repair step as an experiment; the final run 7 uses only the settings
  required by the handout.}
  \label{tab:partD_runs}
\end{table}

Runs 1 to 3 show slack rising by exactly 0.5\,ns per step, so the slow-corner
path was a fixed $\approx 5.25$\,ns. Runs 4 and 5 did not follow this pattern:
with a looser clock, the area-oriented flow chose smaller, slower cells (for
example a \texttt{mux2\_1} of 0.80\,ns and \texttt{\_2} instead of \texttt{\_4}
drive strengths on the critical path), and the slack ended just below zero
both times. At 6.5\,ns the netlist was the same as at 6.0\,ns, so the extra
0.5\,ns appeared fully as slack. Run 6 showed that the 6.0\,ns design can also
pass if an extra timing repair step is enabled after routing.

\subsection{Physical Design Results}

\begin{table}[H]
  \centering
  \begin{tabular}{ll}
    \toprule
    \textbf{Metric} & \textbf{Value} \\
    \midrule
    CLOCK\_PERIOD used & 6.5\,ns \\
    Number of runs needed & 6 (4.0 to 6.5\,ns in 0.5\,ns steps), plus 1 experiment \\
    Core area & 17\,759.5\,$\mu$m$^2$ (die: 22\,500\,$\mu$m$^2$) \\
    Cells after synthesis (Part B) & 68 \\
    Cells after place and route & 4822 \\
    Kinds of cells added & decap, fill, tap, I/O buffers, fanout buffers, \\
     & clock buffers, hold buffer (Table~\ref{tab:partD_cells}) \\
    Wirelength & 2612\,$\mu$m \\
    Setup slack (worst corner, SS 100\,\textdegree C 1.60\,V) & $+0.499$\,ns \\
    Post-route $F_{\max}$ & 166.65\,MHz \\
    $F_{\max}$ from Part C & 336.70\,MHz \\
    DRC errors (Magic / KLayout / router) & 0 / 0 / 0 \\
    LVS errors & 0 \\
    \bottomrule
  \end{tabular}
  \caption{Part D results (final run 7)}
  \label{tab:partD_results}
\end{table}

Hold timing also passes at every corner (worst hold slack $+0.119$\,ns).
The flow reports 11 max-slew warnings, all at the slow corner and all on the
\texttt{rst\_n} reset net, where the fanout buffers drive the reset pins of
the flip-flops. These are not part of the signoff criteria and do not affect
the multiplier data path.

\subsection{Post-Route Maximum Frequency}

\[
  F_{\max,\text{post-route}} = \frac{1000}{\text{CLOCK\_PERIOD} - \text{slack}}
  = \frac{1000}{6.5 - 0.499} = \frac{1000}{6.001} = 166.65\ \text{MHz}
\]

This is about half of the 336.70\,MHz from Part C. Most of the difference
comes from the corner: the same formula with the typical-corner slack
(3.429\,ns) gives $1000 / (6.5 - 3.429) = 325.6$\,MHz, only 3.3\,\% below
Part C. The reasons are discussed in Q4 below.

\subsection{Cells Added by Place and Route}

\begin{table}[H]
  \centering
  \begin{tabular}{llr}
    \toprule
    \textbf{Kind} & \textbf{Cell(s)} & \textbf{Count} \\
    \midrule
    Logic (from synthesis) & \texttt{a21o}, \texttt{xnor2}, \texttt{nand2}, \ldots & 52 \\
    Flip-flops (from synthesis) & \texttt{dfrtp\_2} & 16 \\
    \midrule
    Decoupling capacitors & \texttt{decap\_3} & 4214 \\
    Filler cells & \texttt{fill\_1} (211), \texttt{fill\_2} (61) & 272 \\
    Tap cells & \texttt{tapvpwrvgnd\_1} & 245 \\
    Input / output buffers & \texttt{clkdlybuf4s25\_1} (9 in, 8 out) & 17 \\
    Fanout buffers (\texttt{rst\_n}) & \texttt{clkdlybuf4s25\_1} & 2 \\
    Clock tree buffers & \texttt{clkbuf\_16} (1 root, 2 leaf) & 3 \\
    Hold-fix buffer & \texttt{dlygate4sd3\_1} & 1 \\
    \midrule
    \textbf{Total} & & \textbf{4822} \\
    \bottomrule
  \end{tabular}
  \caption{Cell count after place and route, by kind}
  \label{tab:partD_cells}
\end{table}

The 68 cells from synthesis are all still present (52 logic cells and 16
flip-flops), now at drive strength \texttt{\_2}. Place and route added
$4822 - 68 = 4754$ cells. Of these, 4731 (decap, fill and tap) are physical
cells with no logic function, and only 23 are buffers. The functional cells
cover $1525\,\mu$m$^2$, only 8.6\,\% of the core; decap and fill cells fill
the remaining space.

\subsection{Layout Screenshots}

\begin{figure}[H]
    \centering
    \includegraphics[width=0.8\linewidth,height=0.45\textheight,keepaspectratio]{layout_full_die.png}
    \caption{Final layout of the whole $150\,\mu$m $\times$ $150\,\mu$m die
    in KLayout. The core is covered by standard cell rows and the power grid
    straps; the I/O pins are on the die edges.}
    \label{fig:layout_full}
\end{figure}

\begin{figure}[H]
    \centering
    \includegraphics[width=\linewidth,height=0.4\textheight,keepaspectratio]{layout_zoom.png}
    \caption{Zoomed view near the centre of the core: horizontal standard cell
    rows separated by the VPWR/VGND rails, logic cells between decap and fill
    cells, and horizontal and vertical metal signal wires joined by vias.}
    \label{fig:layout_zoom}
\end{figure}

\subsection{Report Questions}

\subsubsection{Q4: Post-Route $F_{\max}$ and Added Cells}

The post-route $F_{\max}$ is 166.65\,MHz, compared with 336.70\,MHz in
Part C. It is lower for these reasons:
\begin{itemize}
  \item \textbf{Slow corner.} Part C used only the typical corner (TT,
        25\,\textdegree C, 1.80\,V). LibreLane signs off at the worst corner
        (SS, 100\,\textdegree C, 1.60\,V), where every cell is much slower.
        This is the largest effect: at the typical corner the post-route
        $F_{\max}$ would be 325.6\,MHz.
  \item \textbf{Real wires.} Part C assumed ideal wires. After routing, the
        2612\,$\mu$m of wire is extracted, and its resistance and capacitance
        add delay to every cell.
  \item \textbf{Real clock.} Part C used an ideal clock. After clock tree
        synthesis the clock arrives at each flip-flop through buffers with
        some skew, and the flow adds a 0.25\,ns clock uncertainty margin.
  \item \textbf{Area-oriented cell choice.} With a loose period, the flow
        uses smaller, slower cells, since it only needs to meet the target
        period (seen in runs 4 and 5).
\end{itemize}

The flow added these kinds of cells:
\begin{itemize}
  \item \textbf{Tap cells} (245, \texttt{tapvpwrvgnd}) connect the n-well to
        VPWR and the substrate to VGND at regular intervals. Without them the
        wells could float and the chip could suffer latch-up.
  \item \textbf{Decap cells} (4214, \texttt{decap\_3}) are small capacitors
        between VPWR and VGND. They supply charge locally when many cells
        switch at once, which reduces supply noise (IR drop and ground
        bounce).
  \item \textbf{Filler cells} (272, \texttt{fill\_1}, \texttt{fill\_2}) fill
        the remaining gaps in each row so that the n-well, implant layers
        and power rails are continuous, which the DRC rules require.
  \item \textbf{Clock buffers} (3, \texttt{clkbuf\_16}) form the clock tree
        so that the clock reaches all 16 flip-flops at nearly the same time
        with sharp edges.
  \item \textbf{Buffers} (20): input and output buffers isolate the internal
        logic from the external pins, fanout buffers let the \texttt{rst\_n}
        signal drive all 16 flip-flop resets, and one delay buffer
        (\texttt{dlygate4sd3}) slows down a short path to fix a hold
        violation.
\end{itemize}

\subsubsection{Q5: DRC and LVS}

\textbf{DRC} (design rule check) verifies that the layout geometry follows
the foundry's manufacturing rules: minimum widths and spacings of every
layer, minimum enclosure of vias and contacts, minimum area, and density.
\textbf{LVS} (layout versus schematic) extracts the transistors and
connections from the layout and checks that they form exactly the same
circuit as the gate-level netlist: same devices, same nets and same pins.

Both must pass before manufacture because they catch different failures.
A DRC error means the layout may not be built as drawn: wires that are too
thin or too close can break or short, which lowers yield or kills the chip.
An LVS error means the layout can be built correctly but is the wrong
circuit, for example a missing connection or two shorted nets, so the chip
would not compute $a \times b$. Our final layout has 0 DRC errors in Magic,
KLayout and the router, and 0 LVS errors with no unmatched devices, nets or
pins, so the layout is both manufacturable and matches the netlist.


\end{document}
